\section{Kondo 问题}

Gross 第 34 章用费米液体语言讨论磁杂质：电阻极小、Kondo 温度，以及 $T\ll T_K$ 的局域费米液体。微观模型见第 2 章 Anderson/Kondo 哈密顿量。

\subsection{现象}

纯金属低温 $\rho\sim T^5$（电声）。掺入稀磁杂质（Fe、Mn 等）后，$\rho(T)$ 出现极小；$T_{\min}$ 随杂质浓度 $n_i$ 变化，极小深度 $\propto n_i$。这超出简单势能散射。

\subsection{Kondo 的对数}

交换
\begin{equation}
  H_K
  =\sum_{\mathbf{k}\sigma}\xi_{\mathbf{k}}c_{\mathbf{k}\sigma}^\dagger c_{\mathbf{k}\sigma}
  +J\,\mathbf{S}\cdot\mathbf{s}(0),
\end{equation}
$J>0$ 为反铁磁。二阶自旋翻转过程因费米海 Pauli 阻塞产生
\begin{equation}
  J_{\mathrm{eff}}(\omega)
  =J+J^2 N(0)\ln(D/|\omega|)+\cdots.
\end{equation}
电阻
\begin{equation}
  \rho(T)
  \approx
  \rho_0
  \bigl[1+2J N(0)\ln(D/T)+\cdots\bigr]^2
\end{equation}
随降温上升，与电声 $T^5$ 相加即给出极小。对数在
\begin{equation}
  T_K\sim D\exp\bigl(-1/JN(0)\bigr)
\end{equation}
处发散，微扰崩溃。

\subsection{低能费米液体}

$T\ll T_K$ 时杂质自旋被传导电子单态屏蔽，固定点是\textbf{局域费米液体}（Nozi\`eres）：
\begin{itemize}
  \item 费米面相移 $\delta=\pi/2$（幺正极限），$s$ 波共振；
  \item 比热、磁化率由同一 Wilson 比 $R=\pi^2 k_B^2\chi/(3\mu_B^2\gamma)=2$ 联系；
  \item $\Delta\rho(T)\propto (T/T_K)^2$。
\end{itemize}
Anderson 模型在 $U\to\infty$、深能级极限经 Schrieffer--Wolff 变换化为 Kondo。完整温度标度需数值重整化群；费米液体只控制红外。

\begin{note}
Gross 强调：即便强关联杂质，低能仍可能用少数 Landau 参数描述。这与第 7 章“准粒子是普适红外语言”的主题一致。
\end{note}
